\subsection{Summary of Chapter 2}

This chapter translated the physical and geological ideas introduced in Chapter 1 into a mathematical formulation of the thermoelastic wave propagation problem. The geological medium was represented as a continuum domain \(\Omega \subset \mathbb{R}^3\), with position \(x=(x, y)\) and time \(t\in[0,T_f]\). The main physical fields were introduced, including temperature, temperature perturbation, displacement, velocity, acceleration, strain, and stress.

The chapter also defined the material parameters required for the thermoelastic formulation. These include density, elastic moduli, thermal expansion coefficient, thermal conductivity, heat capacity, volumetric heat capacity, and the thermoelastic coupling coefficient. In the simplified mathematical setting, these quantities are treated as effective parameters of a locally homogeneous isotropic linear thermoelastic medium. At the same time, the formulation recognizes that real geological materials may be heterogeneous, porous, fractured, pressure-dependent, temperature-dependent, and anisotropic.

The governing equations were then assembled into a coupled system. The equation of motion describes acceleration caused by stress gradients. The strain--displacement relation defines deformation from displacement. The thermoelastic constitutive relation connects stress with mechanical strain and temperature perturbation. The coupled heat equation describes heat conduction and the interaction between temperature evolution and volumetric deformation. Together, these equations define the thermal--mechanical structure of the problem.

The chapter also introduced the initial and boundary conditions required for a complete mathematical problem. Initial displacement, initial velocity, and initial temperature define the starting state of the medium, while mechanical and thermal boundary conditions describe how the selected rock domain interacts with its surroundings. Finally, directional propagation quantities were defined using source and observation points, the propagation vector, the unit direction vector, distance, azimuth, and elevation.

The formulation is written in general two-dimensional form because real geological media are spatially two-dimensional. However, the thesis does not claim to provide a complete high-fidelity two-dimensional field-scale solution of the full coupled thermoelastic system. Instead, the formulation establishes the physical variables, governing relations, assumptions, and directional quantities needed for later interpretation. In this way, Chapter 2 provides the mathematical bridge between the physical foundations of Chapter 1 and the subsequent analysis of directional thermoelastic wave behavior.